% ============================================================================
% Section VI: Validation and Benchmark Results
% ============================================================================

\section{Validation and Benchmark Results}
\label{sec:validation}

\gtopt{} is validated on IEEE benchmark networks and derived cases,
on published SDDP, unit-commitment, reserve and gas-network benchmarks,
and against independent models of the same problems.  All IEEE test
cases converge to optimal solutions (status~0) with zero optimality
gap; Table~\ref{tab:benchmarks} summarizes them.

\begin{table*}[t]
  \centering
  \caption{IEEE Benchmark Results (scale\_objective = 1000, CLP solver)}
  \label{tab:benchmarks}
  \begin{tabular}{@{}llcccccrc@{}}
    \toprule
    \textbf{Case} & \textbf{Type} & \textbf{Buses} & \textbf{Gen}
      & \textbf{Dem} & \textbf{Lines} & \textbf{Bat}
      & \textbf{Obj (scaled)} & \textbf{Blocks} \\
    \midrule
    \multicolumn{9}{@{}l}{\textit{Standard IEEE DC OPF (single snapshot)}} \\
    \texttt{ieee\_4b\_ori}   & DC OPF     & 4  & 2 & 2  & 5  & -- & 5.000    & 1  \\
    \texttt{ieee\_9b\_ori}   & DC OPF     & 9  & 3 & 3  & 9  & -- & 55.184   & 1  \\
    \texttt{ieee\_14b\_ori}  & DC OPF     & 14 & 5 & 11 & 20 & -- & 8.334    & 1  \\
    \texttt{ieee30b}          & DC OPF     & 30 & 6 & 21 & 41 & -- & 5.668    & 1  \\
    \texttt{ieee\_57b}       & DC OPF     & 57 & 7 & 42 & 80 & -- & 25.016   & 1  \\
    \midrule
    \multicolumn{9}{@{}l}{\textit{24-hour dispatch with solar profiles}} \\
    \texttt{ieee\_9b}        & Solar+OPF  & 9  & 3 & 3  & 9  & -- & 5400.007 & 24 \\
    \texttt{ieee\_14b}       & Solar+OPF  & 14 & 5 & 11 & 20 & -- & 1492.283 & 24 \\
    \midrule
    \multicolumn{9}{@{}l}{\textit{Battery energy storage}} \\
    \texttt{bat\_4b}         & BESS+OPF   & 4  & 3 & 2  & 5  & 1  & 14.160   & 4  \\
    \texttt{bat\_4b\_24}     & BESS+OPF   & 4  & 3 & 2  & 5  & 1  & 44.862   & 24 \\
    \midrule
    \multicolumn{9}{@{}l}{\textit{Capacity expansion}} \\
    \texttt{exp\_gen\_4b}    & Gen. exp.  & 4  & 2 & 2  & 5  & -- & 5.029    & 4  \\
    \texttt{exp\_bat\_4b}    & Bat. exp.  & 4  & 3 & 2  & 5  & 1  & 12.437   & 4  \\
    \texttt{c0}              & Dem. exp.  & 1  & 1 & 1  & 0  & -- & 23.163   & 5 stages \\
    \bottomrule
  \end{tabular}
\end{table*}

\subsection{Test Systems and Data Sources}
\label{sec:test_systems}

Table~\ref{tab:test_systems} lists the public test systems included in
the repository (directory \texttt{cases/}) or read by its converters,
with their sources.  Each case in the repository carries reference
results that are checked in continuous integration.  The standard test
systems, the example systems of other tools and the published case
studies are also distributed, converted with the converters of
\gtopt{}, in the companion case repository~\cite{gtoptcases}, whose
benchmark collection keeps the upstream files of openly licensed data
and, for data without a licence, a recipe that downloads a pinned,
checksum-verified copy and converts it locally.

\begin{table*}[t]
  \centering
  \caption{Public test systems used for validation and their sources}
  \label{tab:test_systems}
  \small
  \setlength{\tabcolsep}{4pt}
  \begin{tabular}{@{}llll@{}}
    \toprule
    \textbf{System} & \textbf{Size} & \textbf{Used for} & \textbf{Source} \\
    \midrule
    IEEE 4-bus and battery variants & 4 buses & DC OPF, storage, expansion & illustrative \\
    Anderson--Fouad 9-bus & 9 buses & DC OPF, solar profile & \cite{anderson2002power} \\
    IEEE 14-, 30-, 57-bus & 14--57 buses & DC OPF, pandapower cross-check & \cite{pstca, ieee30bus, ieee57bus, thurner2018pandapower} \\
    IEEE 118-bus & 118 buses & DC OPF & \cite{pstca, pglibopf2019} \\
    PGLib-OPF v23.07 (5 to 300 buses) & 5--300 buses & DC OPF & \cite{pglibopf2019} \\
    IEEE RTS-96 (24-bus) & 24 buses & unit commitment with OPF & \cite{grigg1999rts96, pglibopf2019} \\
    RTS-GMLC & 73 buses & UC and dispatch with renewables & \cite{barrows2020rtsgmlc} \\
    NREL-118 & 118 buses & one-week hourly dispatch & \cite{pena2018nrel118, powersimulations} \\
    NREL Sienna test systems & 1--5 buses & UC, hydro, storage, reserves & \cite{lara2021powersystems, powersimulations} \\
    UnitCommitment.jl fixtures & 14--118 buses & UC MIP vs.\ UnitCommitment.jl & \cite{xavier2024ucjl} \\
    Clustered-UC case studies & 39, 118 buses & UC MIP vs.\ published model & \cite{morales2019cuc, cucdata} \\
    Power-based UC example & 1 bus & published schedule and cost & \cite{philipsen2019trading} \\
    Chebyshev scheduling systems & 1--9 buses & block vs.\ point-sampled stages & \cite{matus2015chebyshev} \\
    Non-convex pricing examples & 1--3 buses & UC prices, uplift & \cite{oneill2005pricing, gribik2007convexhull, hua2017convexhull} \\
    MATPOWER \texttt{case30} reserves & 30 buses & zonal reserves, interfaces & \cite{matpower_case30} \\
    Badesa et al.\ frequency services & 1 bus & reserve and inertia prices & \cite{badesa2020pricing} \\
    Polish \texttt{case3375wp} (R.~Korab) & 3\,375 buses & large-network DC OPF & \cite{zimmerman2011matpower} \\
    SDDP.jl examples (first steps, hydro valleys, FAST) & 1--2 reservoirs & SDDP vs.\ published bounds & \cite{dowson2020sddp_jl, cambier_fast} \\
    PSR SDDP sample cases & multi-system & \texttt{sddp2gtopt} hydrothermal & \cite{psrclasses} \\
    Saldarriaga pipe chain & 4 gas nodes & Weymouth pressures & \cite{saldarriaga2018phd} \\
    GasLib-11 & 11 gas nodes & radial gas network & \cite{schmidt2017gaslib, pfetsch2015nominations} \\
    Belgian network (GasModels \texttt{A1}) & 24 gas nodes & radial gas network & \cite{dewolf2000gas, gasmodels} \\
    GasModels \texttt{case-6} & 6 gas nodes & exact Weymouth, meshed SLP & \cite{gasmodels} \\
    Hu--Ryan 6-bus/7-node, IEEE-24/Belgian-20 & 6--24 buses, 7--20 gas nodes & coupled gas--power & \cite{hu2019data, hu2019stochastic} \\
    Ratha et al.\ 24-bus/12-node & 24 buses, 12 gas nodes & meshed coupled gas--power & \cite{ratha2020data, ratha2020affine, ordoudis2019gas} \\
    CEN daily schedules (PLEXOS) & $\sim$250 buses & vs.\ official schedule & \cite{cen_pcp}, Section~\ref{sec:cases} \\
    Real systems of six countries & 3--3\,374 buses & UC, SDDP, OPF & Section~\ref{sec:gtopt_cases} \\
    \bottomrule
  \end{tabular}
\end{table*}

\subsection{Public Real-System Cases}
\label{sec:gtopt_cases}

Beyond the benchmark networks of Table~\ref{tab:test_systems}, GTOpt
is exercised on real power systems through the public companion
repository \texttt{gtopt-cases}~\cite{gtoptcases2026}.  Each case is
built with the GTOpt converters from data published by a system
operator, market operator or regulator, and is distributed with its
source data, the converted case, the exact command to run and rebuild
it, and the provenance (converter revision, flags and file hashes).
Only inputs are published, so that results can be reproduced and
compared independently; the cases are provided as is and are not
validated by the publishing institutions, which keep all rights over
the data.  Table~\ref{tab:gtopt_cases} lists the six countries
covered.  The Irish data may not be redistributed, so that case is
published as a recipe that downloads the public files and converts
them locally.  The 23-year Chilean cases are committed for their first
year to limit the repository size; the full horizon is built from the
same sources with one converter command.

\begin{table*}[t]
  \centering
  \caption{Real-system cases in \texttt{gtopt-cases}}
  \label{tab:gtopt_cases}
  \small
  \begin{tabular}{@{}lllll@{}}
    \toprule
    \textbf{Country} & \textbf{System · publisher} & \textbf{Cases} & \textbf{Horizon} & \textbf{Converter} \\
    \midrule
    Australia & NEM · AEMO & 1 dispatch, 2026 ISP model (Step Change) & 1 week, hourly & \texttt{plexos2gtopt} \\
    Chile & SEN · CEN & 14 daily operation programmes (unit commitment) & 1 week each & \texttt{plexos2gtopt} \\
          &            & 3 hydrothermal SDDP (PLP), irrigation agreements & 2, 4.5 and 23 years & \texttt{plp2gtopt} \\
          &            & 3 PET 2027 transmission planning bases (PLP) & 23 years & \texttt{plp2gtopt} \\
    Guatemala & SNI · AMM, EOR & 2 NCP daily/weekly dispatches & 1 week & \texttt{sddp2gtopt} \\
              &                & 6 PSS\textregistered E network snapshots & 1 hour & \texttt{psse2gtopt} \\
              &                & 6 annual SDDP programmes & 1 year & \texttt{sddp2gtopt} \\
    Ireland & SEM · SEM Committee (CRU, UR) & 1 unit commitment (recipe only) & 1 week, hourly & \texttt{plexos2gtopt} \\
    Peru & SEIN · COES & 4 SDDP cases, 2026--2029 reserve-margin study & 1 year each & \texttt{sddp2gtopt} \\
    Poland & MATPOWER \texttt{case3375wp} & 1 DC OPF, 3\,374 buses & 1 hour & \texttt{pp2gtopt} \\
    \bottomrule
  \end{tabular}
\end{table*}

\subsection{IEEE 4-Bus Network}
\label{sec:ieee4b}

The IEEE 4-bus test case consists of a fully meshed network with 4 buses
and 5 transmission lines (Fig.~\ref{fig:ieee4b}).  Two generators
(g1: 300~MW at \$20/MWh, g2: 200~MW at \$35/MWh) supply two demands
(d3: 150~MW, d4: 100~MW).

\begin{figure}[t]
  \centering
  \begin{tikzpicture}[
    >=Stealth,
    bus/.style={circle, draw, thick, fill=white, minimum size=0.7cm,
      font=\scriptsize\bfseries},
    genbox/.style={circle, draw, thick, fill=yellow!20,
      minimum size=0.6cm, font=\scriptsize},
    dembox/.style={-{Triangle[length=5pt,width=4pt]}, thick, red!70},
    line/.style={thick, blue!60},
    linelab/.style={font=\tiny, fill=white, inner sep=1pt},
  ]
    % Buses in diamond layout
    \node[bus] (b1) at (0,2)   {$b_1$};
    \node[bus] (b2) at (3,2)   {$b_2$};
    \node[bus] (b3) at (0,0)   {$b_3$};
    \node[bus] (b4) at (3,0)   {$b_4$};

    % Generators
    \node[genbox] (g1) at (-1.3,2) {$\sim$};
    \node[font=\tiny, above=1pt of g1] {g1: 300\,MW};
    \node[font=\tiny, below=1pt of g1] {\$20/MWh};
    \draw[->, thick] (g1) -- (b1);

    \node[genbox] (g2) at (4.3,2) {$\sim$};
    \node[font=\tiny, above=1pt of g2] {g2: 200\,MW};
    \node[font=\tiny, below=1pt of g2] {\$35/MWh};
    \draw[->, thick] (g2) -- (b2);

    % Demands (downward arrows from buses)
    \draw[dembox] (b3) -- ++(0,-0.8)
      node[below, font=\tiny] {d3: 150\,MW};
    \draw[dembox] (b4) -- ++(0,-0.8)
      node[below, font=\tiny] {d4: 100\,MW};

    % Transmission lines
    \draw[line] (b1) -- (b2) node[linelab, midway, above] {$\ell_{12}$};
    \draw[line] (b1) -- (b3) node[linelab, midway, left]  {$\ell_{13}$};
    \draw[line] (b2) -- (b3) node[linelab, pos=0.5, below right=2pt and 0pt]
      {$\ell_{23}$};
    \draw[line] (b2) -- (b4) node[linelab, midway, right] {$\ell_{24}$};
    \draw[line] (b3) -- (b4) node[linelab, midway, below] {$\ell_{34}$};
  \end{tikzpicture}
  \caption{IEEE 4-bus network: 2 generators, 2 demands, 5 lines forming
    a fully meshed topology for DC OPF validation.}
  \label{fig:ieee4b}
\end{figure}

Under DC OPF with Kirchhoff constraints, the cheaper generator g1
dispatches 250~MW and g2 idles.  The scaled objective value is
$5.0 = 250 \times 20 / 1000$.  Power routes through parallel paths
(b1$\rightarrow$b3 and b1$\rightarrow$b2$\rightarrow$b3) according to
the relative reactances.

\subsection{IEEE 9-Bus Network}
\label{sec:ieee9b}

The Anderson--Fouad 9-bus system~\cite{anderson2002power} has three
generation areas connected by a meshed transmission network
(Fig.~\ref{fig:ieee9b}).  Generators g1 (250~MW, \$20/MWh),
g2 (300~MW, \$35/MWh), and g3 (270~MW, \$30/MWh) serve demands at
buses 5, 7, and~9 (total 315~MW).

\begin{figure}[t]
  \centering
  \begin{tikzpicture}[
    >=Stealth,
    bus/.style={circle, draw, thick, fill=white, minimum size=0.55cm,
      font=\tiny\bfseries, inner sep=1pt},
    genbus/.style={bus, fill=yellow!15},
    demlab/.style={font=\tiny, red!70},
    genlab/.style={font=\tiny, green!40!black},
    line/.style={thick, blue!50},
    dem/.style={-{Triangle[length=3pt,width=3pt]}, thick, red!60},
  ]
    % network ring 4-5-7-8-9-6-4, generator buses 1, 2, 3 outside
    \node[bus] (b4) at (0,1.9)     {4};
    \node[bus] (b5) at (-1.5,1.0)  {5};
    \node[bus] (b6) at (1.5,1.0)   {6};
    \node[bus] (b7) at (-1.5,-0.5) {7};
    \node[bus] (b9) at (1.5,-0.5)  {9};
    \node[bus] (b8) at (0,-1.4)    {8};
    \node[genbus] (b1) at (0,3.1)    {1};
    \node[genbus] (b2) at (-2.9,-0.5) {2};
    \node[genbus] (b3) at (2.9,-0.5)  {3};

    \foreach \a/\b in {1/4, 4/5, 4/6, 5/7, 6/9, 7/8, 8/9, 2/7, 3/9}
      \draw[line] (b\a) -- (b\b);

    \node[genlab, right=3pt of b1] {g1: 250\,MW, \$20};
    \node[genlab, above=2pt of b2] {g2: 300\,MW, \$35};
    \node[genlab, above=2pt of b3] {g3: 270\,MW, \$30};

    \draw[dem] (b5.west) -- ++(-0.5,0) node[demlab, left] {125\,MW};
    \draw[dem] (b7.south) -- ++(0,-0.45) node[demlab, below] {100\,MW};
    \draw[dem] (b9.south) -- ++(0,-0.45) node[demlab, below] {90\,MW};
  \end{tikzpicture}
  \caption{IEEE 9-bus network: three generation areas (buses 1--3)
    serve three demands (buses 5, 7, 9) through a meshed
    9-line transmission network.}
  \label{fig:ieee9b}
\end{figure}

In the single-snapshot case (\texttt{ieee\_9b\_ori}), economic dispatch
follows merit order constrained by network topology.  The 24-hour variant
(\texttt{ieee\_9b}) replaces g3 with a solar plant (\$0/MWh) whose output
follows a daily capacity profile peaking at noon.  During daylight hours,
zero-cost solar generation displaces g2; at night, g2 carries the load.

\subsection{IEEE 14-Bus Network}
\label{sec:ieee14b}

The 14-bus system has 5 generators, 11 demands, and 20 transmission lines.
The constrained variant (\texttt{ieee\_14b}) tightens limits on lines
$\ell_{1-2}$ and $\ell_{1-5}$, forcing both to saturate simultaneously
at the evening demand peak.  This validates that shadow prices (Lagrange
multipliers of the binding Kirchhoff constraints) are computed correctly.

\subsection{IEEE 30-Bus and 57-Bus Networks}

The 30-bus~\cite{ieee30bus} and 57-bus~\cite{ieee57bus} systems test
scalability.  The 30-bus case has 6 generators, 21 demands, and 41 lines;
the 57-bus case has 7 generators, 42 demands, and 80 lines.  Both
converge to optimal solutions in under one second with the CLP solver.

\subsection{Battery Storage Cases}
\label{sec:bat_cases}

The \texttt{bat\_4b} case augments the 4-bus network with a 200~MWh
battery at bus~3 (60~MW charge/discharge, 95\% round-trip efficiency) and
a solar generator with a 4-period capacity profile.  The battery charges
during solar peak and discharges during evening demand, reducing total
system cost compared to the no-storage baseline.

The 24-hour variant (\texttt{bat\_4b\_24}) demonstrates full daily
charge/discharge cycling with state-of-charge tracking across all
24~blocks.

\subsection{Capacity Expansion Cases}
\label{sec:exp_cases}

Three expansion cases validate the investment planning formulation:

\begin{itemize}
  \item \texttt{exp\_gen\_4b}: Generator g1 can be expanded by up to
    290~MW at \$1000/year annualized cost.  The optimizer adds capacity
    to avoid dispatching expensive g2.
  \item \texttt{exp\_bat\_4b}: Battery expansion with 200~MWh modules
    at \$1000/year.  The optimizer installs storage to arbitrage
    solar/thermal cost differentials.
  \item \texttt{c0}: Multi-stage demand expansion (5 stages, 10\%
    discount rate).  Demand grows through 20~MW modules at
    \$8760/year.  The optimizer balances discounted investment cost
    against increasing served demand.
\end{itemize}

\subsection{Cross-Validation with pandapower}

The \texttt{gtopt\_compare} tool performs automated cross-validation
against pandapower~\cite{thurner2018pandapower} DC OPF on all IEEE test
networks (4, 9, 14, 30, 57 buses).  For each case, both tools solve the
same network data and compare:

\begin{itemize}
  \item Objective function value (total generation cost)
  \item Generator dispatch (MW per generator)
  \item Bus angles (radians)
  \item Line flows (MW)
  \item Locational marginal prices (bus balance duals)
\end{itemize}

Results match to within solver tolerance ($<10^{-6}$ relative error)
for all standard IEEE cases, confirming the correctness of the DC OPF
implementation in \gtopt{}.

\subsection{SDDP Against Published Benchmarks}
\label{sec:sddp_bench}

The SDDP implementation is checked in continuous integration against
published reference values:
\begin{itemize}
  \item \textbf{SDDP.jl \emph{first steps}}~\cite{dowson2020sddp_jl}: a
    three-stage hydrothermal problem with stagewise-independent inflows
    and a published lower bound of 8333.33, confirmed in-repo by the
    27-leaf extensive form.  Every iterate of \gtopt{} is a valid lower
    bound and the sequence is monotone.  With the default aperture
    backward pass it converges to 8055.6, 3.3\% below the exact value,
    because the cuts are built from the resampled apertures rather than
    from the exact expectation.  Depending on the dual vertex returned
    by the LP solver, it can also reach the exact value.
  \item \textbf{SDDP.jl \emph{Markov uncertainty}} and
    \emph{generation expansion}: checked against extensive-form and
    MIP oracles.
  \item \textbf{Infanger--Morton dual sharing}~\cite{infanger1996cutsharing}:
    cut-for-cut agreement
    of the shared-dual backward pass with the standard one where the
    theory guarantees it.
\end{itemize}

In addition, the equivalence between the monolithic LP and SDDP is
tested on multi-stage cases: with the same data, the converged SDDP
bounds match the monolithic optimum within the convergence tolerance.
The water value of a spilling reservoir is also checked against
hand-derived values, with free and priced exits, in monolithic and
SDDP runs.

\subsection{Unit-Commitment and Reserve Benchmarks}
\label{sec:uc_bench}

The commitment, reserve and pricing models of
Section~\ref{sec:formulation} are cross-checked against other
implementations and against published values:
\begin{itemize}
  \item \textbf{UnitCommitment.jl}~\cite{xavier2024ucjl}: its IEEE
    14-bus and 118-bus test fixtures (tiered startup costs, minimum
    up/down times, ramps, initial conditions, profiled renewables) are
    converted to \gtopt{} and solved as MIPs; the hourly commitment
    and dispatch are compared with the solution of UnitCommitment.jl
    on the same instance, with the full network and with a copper
    plate.
  \item \textbf{NREL Sienna}~\cite{lara2021powersystems,
    powersimulations}: the PowerSimulations.jl test systems of the
    5-bus PJM network (cascaded hydro, pumped storage, an HVDC link, an
    interruptible load, a single monitored line, time-varying fuel
    prices, a battery) are compared on hourly and aggregate dispatch,
    and its reserve systems (range reserves and operating-reserve
    demand curves, represented as stacked priced reserve zones) on the
    objective to a relative $10^{-6}$ and on the hourly reserve totals.
  \item \textbf{Clustered unit commitment}~\cite{morales2019cuc}: the
    modified IEEE 39-bus and 118-bus systems of Morales-Espa\~na and
    Tejada-Arango, built from the published data~\cite{cucdata} as
    individual-unit commitment problems.  An independent transcription
    of the published model reaches the same incumbents as \gtopt{} to
    the cent.  The objectives reported in the paper are higher by
    $0.8\%$ to $2.6\%$ over its four cases, by an amount that grows
    with the reserve requirement and, for the 39-bus system, vanishes
    when extrapolated to zero reserve; this points to a different
    reserve representation in the code behind the published table.  The
    block formulation of Section~\ref{sec:sc}, applied to the same hours
    as one non-chronological stage, is checked to remain a relaxation of
    the chronological commitment.
  \item \textbf{Power-based commitment}: the day-ahead example of
    Philipsen et al.~\cite{philipsen2019trading} (two ramp-limited
    units, seven hours) is solved with the formulation of
    Section~\ref{sec:pbuc}; the published power-based cost and the
    boundary powers of both units are reproduced exactly, and the
    energy-block formulation reproduces the lower energy-based cost of
    the same paper, which the ramp limits cannot actually deliver.
  \item \textbf{Point-sampled stages}: small hydrothermal, battery,
    network, commitment and periodic cases, each in an hourly-block and
    a Chebyshev--Lobatto point version, are converted from the
    validation systems of the Chebyshev scheduling
    code~\cite{matus2015chebyshev}; their objectives are checked
    against hand-derived or closed-form values and as regressions, and
    the storage of every rule is checked to be conserved over a closed
    period.
  \item \textbf{Non-convex pricing}: the fixed, relaxed and fast-start
    pricing modes reproduce the prices and uplifts of the published
    examples of O'Neill et al.~\cite{oneill2005pricing}, Gribik, Hogan
    and Pope~\cite{gribik2007convexhull} and Hua and
    Baldick~\cite{hua2017convexhull}; the remaining differences are
    explained by dual degeneracy or by ramp-coupled relaxations that are
    not the convex hull.
  \item \textbf{Reserves}: the MATPOWER zonal-reserve
    case~\cite{matpower_case30} (two overlapping zones and interface
    limits) is reproduced exactly, including the zone prices and the
    reserve cost, and so are the economic-dispatch tables of Badesa et
    al.~\cite{badesa2020pricing} for energy and frequency-response
    prices.
  \item \textbf{Emissions}: the coal--gas dispatch of the
    myopic-pathway example of PyPSA~\cite{brown2018pypsa}, reduced to
    one period with a CO$_2$ cap, and a three-zone combined-cycle
    system with the technology parameters of the three-zone example of
    GenX~\cite{jenkins2017genx} are solved with a binding cap; the
    dispatch and the carbon price (the dual of the cap) are checked
    against hand-derived solutions.
\end{itemize}

\subsection{Natural-Gas Network Benchmarks}
\label{sec:gas_bench}

The gas network of Section~\ref{sec:gas} is checked in continuous
integration on the cases of Table~\ref{tab:gas_bench}.  Each reference
comes either from a published value or from an independent model built
outside \gtopt{}: an LP of the same tangent formulation (SciPy with
HiGHS~\cite{huangfu2018highs}) or an exact nonlinear Weymouth solve.
The coupled cases are LP reductions of the published data (unit
commitment relaxed, deterministic wind, gas storage as a capped
supply), so their objectives are not comparable with the published
stochastic results.

\begin{table*}[t]
  \centering
  \caption{Gas network benchmarks: what each case validates and its
    reference (objective in the case's own units)}
  \label{tab:gas_bench}
  \small
  \begin{tabular}{@{}p{3.3cm}p{6.2cm}p{6.6cm}@{}}
    \toprule
    \textbf{Case} & \textbf{Validates} & \textbf{Reference} \\
    \midrule
    Saldarriaga Fig.~2.2(b)~\cite{saldarriaga2018phd}
      & three pipes in series, 32 tangents
      & published node pressures 90/70/50/30~bar within the tangent
        bound (0.05~bar); objective of the independent LP \\
    GasLib-11, valve closed~\cite{schmidt2017gaslib}
      & radial network with compressors
      & independent LP (295.41); exact Weymouth drop on the LP flows \\
    Belgian \texttt{A1}~\cite{dewolf2000gas, gasmodels}
      & 19 pipes and 4 stations after merging parallel arcs
      & independent LP (8849.43); shedding only at the node whose
        pressure floor cannot be met \\
    GasModels \texttt{case-6}~\cite{gasmodels}
      & loop with fixed orientation and 64 tangents; the same loop with
        linearized pipes and SLP
      & published exact optimum $-253.683$: the tangent LP is within
        $1.6\cdot10^{-5}$ (relative); the SLP matches the exact NLP to
        $10^{-5}$ \\
    Hu--Ryan 6-bus/7-node~\cite{hu2019data}
      & coupled radial system, 24~h
      & independent LP; the exact Weymouth optimum is 0.68\% higher
        because an upper pressure bound downstream binds, which the
        exactness report flags on that pipe \\
    Hu--Ryan IEEE-24/Belgian-20~\cite{hu2019data}
      & coupled system, radial and SLP variants
      & radial tangent LP within 0.0099\% of the exact optimum; the SLP
        variant reaches the exact optimum \\
    Ratha et al.\ 24-bus/12-node~\cite{ratha2020data}, system
      of~\cite{ordoudis2019gas, ordoudis2018data}
      & meshed coupled system, gas SLP, compressors as ratio-bounded
        stations
      & independent SLP, unchanged by a polish on the exact NLP, i.e.\
        a local optimum of the exact problem \\
    Elevation chain
      & 0--2000--4000~m chain, term~\eqref{eq:gas_elevation}
      & closed-form outlet pressures (56.747 and 44.673~bar) \\
    Component cases
      & compressor fuel models and drives, line pack, plant pressure
        gate, booster and derate
      & independent LP of each variant: objective to $10^{-5}$
        (relative), pressures to 0.01~bar \\
    \bottomrule
  \end{tabular}
\end{table*}

The unit tests additionally check the Weymouth constant and the
line-pack coefficient, the tangent envelope against the exact curve,
the compressor-fuel planes against analytic derivatives, the
convergence of the SLP on loops, under flow reversal and on degenerate
splits, and the replay of the re-linearised rows when LP snapshots are
compressed.

\paragraph{Stylised Chilean studies.}
The gas elements were also applied to stylised studies of the central
Chilean network, in which LNG is regasified at the coast, piped to the
capital region and delivered to combined-cycle plants whose turbines
need a minimum fuel-gas pressure.  These studies explore the interaction
between the minimum plant pressure, a compressor station on the trunk,
an on-site fuel-gas booster, pressure derating and multi-day line pack,
and they cross-check the native elements against an earlier formulation
of the same LP built from generic user-defined variables and
constraints.  The data are plausible values chosen for the study, not
the parameters of the actual pipelines or plants.  Pipelines that cross
the Andes, as in the north of the country, motivated the elevation
term, which changes the attainable delivery pressure substantially on
such routes.

\subsection{Real Systems}
\label{sec:real_systems}

The companion case repository~\cite{gtoptcases} also publishes the
inputs of real systems converted from the public data of their system
operators and regulators, together with the instructions to run them:
daily schedules, medium- and long-term SDDP cases and a transmission
planning case of the Chilean system~\cite{cen_pcp}; daily, weekly and
long-term cases of Guatemala~\cite{amm_eor}; SDDP cases of
Peru~\cite{coes_mrfo}; production-cost and monthly expansion cases of
the Australian National Electricity Market from the 2026 Integrated
System Plan~\cite{aemo2026isp}; unit-commitment weeks of the Irish
Single Electricity Market~\cite{sem2025plexos}; and the Polish
3\,375-bus network~\cite{zimmerman2011matpower}.  The data remain the
property of their publishers; source files are redistributed only where
their terms allow, and otherwise fetched from the publisher.  These
cases exercise the converters from PLEXOS, PLP, PSR SDDP and PSS/E at
national scale; the comparison of Section~\ref{sec:cases} uses the
Chilean daily schedules.
